\scanpage{7}
\[
A_n=A/(F^{n+1},V^{n+1})
\]
\[
I_n=A_n\quad\text{regardé comme \uline{bimodule} sur }A
\]
\[
\underline{I=\underset{n}{\varinjlim} I_n}\qquad I_n\overset p{\unsure{\mathord{-}|\longrightarrow}}I_{n+1}
\]
\[
D:M\leadsto\operatorname{Hom}(M,I)\unsure{'|\mathord{.}}\quad\text{foncteur dualisant}
\]

\uline{Attention}, à priori un module à g. est transformé en module à droite, et vice versa. Mais on peut utiliser l'involution $s$ de $A$, échangeant $F$ et $V$, pour passer de l'un à l'autre

\deleted{$D(M)$}\quad $A_n$ est fini sur $W_n$

Un autre foncteur dualisant est obtenu en faisant (si $M$ tué par $F^n,V^n$)
\[
D':M\leadsto\operatorname{Hom}_{W_n\text{-modules}}(M,W_n)
\]
et définissant sur $D'(M)$ la structure de $A$-module à droite par ``\unsure{transport.|transport|transpos.}''

\[
D'(\unsure{A_n|A|M})\unsure{\simeq|=}\operatorname{Hom}(A_n,W_n)\quad\text{structure de bimodule}
\]
\deleted{\unsure{Or|On}
\[
\operatorname{Hom}_{W_n}(M,W_n)=\operatorname{Hom}_{A_n}\bigl(M,\underbrace{\operatorname{Hom}_{W_n}(A_n,W_n)}_{\substack{\text{structure}\\\text{de bimodule}\\\text{naturelle}}}\bigr)
\]
Je dis qu'on a un \unsure{isom|isom.}
\[
I_n=A_n\simeq\operatorname{Hom}_{W_n}(A_n,W_n)
\]}
\deleted{$uF=u\circ F$}
% archive sheet number: 6

